\begin{thebibliography}{99}
\footnotesize\itemsep=0pt

\bibitem{contra}
Brzezi\'nski T., Hopf-cyclic homology with contramodule coefficients, in
 Quantum Groups and Noncommutative Spaces, \href{https://doi.org/10.1007/978-3-8348-9831-9_1}{\textit{Aspects Math.}}, Vol.~E41,
 Vieweg + Teubner, Wiesbaden, 2011, 1--8, \href{https://arxiv.org/abs/0806.0389}{arXiv:0806.0389}.

\bibitem{BPO_closed}
Bulacu D., Panaite F., Van~Oystaeyen F., Quantum traces and quantum dimensions
 for quasi-{H}opf algebras, \href{https://doi.org/10.1080/00927879908826809}{\textit{Comm. Algebra}} \textbf{27} (1999),
 6103--6122.

\bibitem{CM98}
Connes A., Moscovici H., Hopf algebras, cyclic cohomology and the transverse
 index theorem, \href{https://doi.org/10.1007/s002200050477}{\textit{Comm. Math. Phys.}} \textbf{198} (1998), 199--246,
 \href{https://arxiv.org/abs/math.DG/9806109}{math.DG/9806109}.

\bibitem{CM99}
Connes A., Moscovici H., Cyclic cohomology and {H}opf algebras, \href{https://doi.org/10.1023/A:1007527510226}{\textit{Lett.
 Math. Phys.}} \textbf{48} (1999), 97--108, \href{https://arxiv.org/abs/math.QA/9904154}{math.QA/9904154}.

\bibitem{Drinfeld}
Drinfeld V.G., Quasi-{H}opf algebras, \textit{Leningrad Math.~J.} \textbf{1}
 (1990), 1419--1457.

\bibitem{ENO}
Etingof P., Nikshych D., Ostrik V., Fusion categories and homotopy theory,
 \href{https://doi.org/10.4171/QT/6}{\textit{Quantum Topol.}} \textbf{1} (2010), 209--273, \href{https://arxiv.org/abs/0909.3140}{arXiv:0909.3140}.

\bibitem{hkrs1}
Hajac P.M., Khalkhali M., Rangipour B., Sommerh\"auser Y., Hopf-cyclic homology
 and cohomology with coefficients, \href{https://doi.org/10.1016/j.crma.2003.11.036}{\textit{C.~R.~Math. Acad. Sci. Paris}}
 \textbf{338} (2004), 667--672, \href{https://arxiv.org/abs/math.KT/0306288}{math.KT/0306288}.

\bibitem{hkrs2}
Hajac P.M., Khalkhali M., Rangipour B., Sommerh\"auser Y., Stable
 anti-{Y}etter--{D}rinfeld modules, \href{https://doi.org/10.1016/j.crma.2003.11.037}{\textit{C.~R.~Math. Acad. Sci. Paris}}
 \textbf{338} (2004), 587--590, \href{https://arxiv.org/abs/math.QA/0405005}{math.QA/0405005}.

\bibitem{hks}
Hassanzadeh M., Khalkhali M., Shapiro I., Monoidal categories, 2-traces, and
 cyclic cohomology, \href{http://dx.doi.org/10.4153/CMB-2018-016-4}{\textit{Canad. Math. Bull.}}, {t}o appear,
 \href{https://arxiv.org/abs/1602.05441}{arXiv:1602.05441}.

\bibitem{js}
Jara P., \c{S}tefan D., Hopf-cyclic homology and relative cyclic homology of
 {H}opf--{G}alois extensions, \href{https://doi.org/10.1017/S0024611506015772}{\textit{Proc. London Math. Soc.}} \textbf{93}
 (2006), 138--174, \href{https://arxiv.org/abs/math.KT/0307099}{math.KT/0307099}.

\bibitem{sk}
Kobyzev I., Shapiro I., A categorical approach to cyclic cohomology of
 quasi-{H}opf algebras and {H}opf algebroids, \href{https://arxiv.org/abs/1803.09194}{arXiv:1803.09194}.

\bibitem{majid_closed}
Majid S., Foundations of quantum group theory, \href{https://doi.org/10.1017/CBO9780511613104}{Cambridge University Press},
 Cambridge, 1995.

\bibitem{majid}
Majid S., Quantum double for quasi-{H}opf algebras, \href{https://doi.org/10.1023/A:1007450123281}{\textit{Lett. Math. Phys.}}
 \textbf{45} (1998), 1--9, \href{https://arxiv.org/abs/q-alg/9701002}{q-alg/9701002}.

\bibitem{pv}
Panaite F., Van~Oystaeyen F., A structure theorem for quasi-{H}opf comodule
 algebras, \href{https://doi.org/10.1090/S0002-9939-07-08712-6}{\textit{Proc. Amer. Math. Soc.}} \textbf{135} (2007), 1669--1677,
 \href{https://arxiv.org/abs/math.QA/0506272}{math.QA/0506272}.

\bibitem{SS_closed}
Sak\'alo\v{s} {\v{S}}., On categories associated to a quasi-{H}opf algebra,
 \href{https://doi.org/10.1080/00927872.2016.1175449}{\textit{Comm. Algebra}} \textbf{45} (2017), 722--748, \href{https://arxiv.org/abs/1402.1393}{arXiv:1402.1393}.

\bibitem{s}
Shapiro I., Some invariance properties of cyclic cohomology with coefficients,
 \href{https://arxiv.org/abs/1611.01425}{arXiv:1611.01425}.

\end{thebibliography}
